\section{Normalization and Schema Refinement}
\label{sec:norm}

% Dependency displays are frequent in this section; the default display skips
% leave large gaps between consecutive one-line dependencies.
\setlength{\abovedisplayskip}{5pt plus 2pt minus 2pt}
\setlength{\belowdisplayskip}{5pt plus 2pt minus 2pt}
\setlength{\abovedisplayshortskip}{2pt plus 2pt}
\setlength{\belowdisplayshortskip}{3pt plus 2pt minus 2pt}

This section explains how the team refines the initial environmental data model
into the normalized relational structure used in the final database. The
process begins with problems identified during Entity--Relationship modelling
and mentor review, then continues through structural changes that resolve
duplicate concepts, unsuitable keys, different observation levels, and
repeated information. The team checks the resulting relations through First
Normal Form, Second Normal Form, Third Normal Form, and Boyce--Codd Normal Form
before presenting the final data model in Section~\ref{sec:dm}.

\subsection{Initial Modelling Approach }
\label{subsec:zeronf}
At the beginning, the team prepares a separate Entity--Relationship diagram
for each collected data file. This helps the team understand the main contents
of the files and identify possible entities and attributes. However, during
mentor review, the team finds that using one complete file as one modelling
unit is not suitable in every case.

Some Excel workbooks contain several sheets, and each sheet may hold a
different type of environmental record. For example, the sheets may differ in
their attributes, possible keys, and reporting periods. If the whole workbook
is treated as one unit, unrelated records may be placed under the same entity.
This also makes it difficult to identify the relationships and dependencies
among the attributes.

The team therefore examines each relevant sheet separately. From each sheet,
the team identifies the possible entities, attributes, candidate keys, and
time level of the records. Similar entities from different sheets are then
compared before they are combined. This sheet-based approach gives the team a
clearer starting point for entity integration and normalization.
\begin{figure}[H]
    \centering

    \begin{tikzpicture}[
        node distance=0.7cm,
        every node/.style={
            draw,
            rounded corners,
            align=center,
            minimum width=6.8cm,
            minimum height=0.75cm,
            font=\small
        },
        arrow/.style={
            -{Stealth[length=2.5mm]},
            thick
        }
    ]

    \node (files) {
        Collected Environmental Data Files
    };

    \node (initial) [below=of files] {
        One ERD Prepared for Each File
    };

    \node (review) [below=of initial] {
        Problems Identified During Mentor Review
    };

    \node (sheets) [below=of review] {
        Relevant Sheets Examined Separately
    };

    \node (entities) [below=of sheets] {
        Entities, Attributes, and Possible Keys Identified
    };

    \node (next) [below=of entities] {
        Entity Integration and Normalization
    };

    \draw[arrow] (files) -- (initial);
    \draw[arrow] (initial) -- (review);
    \draw[arrow] (review) -- (sheets);
    \draw[arrow] (sheets) -- (entities);
    \draw[arrow] (entities) -- (next);

    \end{tikzpicture}

    \caption{Steps followed during the initial entity identification}
    \label{fig:initial-entity-identification}
\end{figure}
Figure~\ref{fig:initial-entity-identification} shows how the team changes from
file-based modelling to sheet-based entity identification. The entities
identified at this stage are reviewed and refined during the later
normalization process.

\subsection{Entity Integration }

\label{subsec:entity-integration}

After identifying the entities from the individual sheets, the team compares
them before preparing the combined data model. The comparison considers the
meaning of each entity, its attributes, possible primary key, reporting period,
and relationship with other entities.

Several sources contain similar environmental information. For example,
temperature, rainfall, humidity, and sunshine records appear in more than one
source. Keeping a separate entity for every source would create repeated
structures in the database. Therefore, entities are combined when they
represent the same type of information and use compatible attributes.

Entities are not combined only because they have similar names. The team also
checks the level of each record. Some sources provide daily observations,
whereas others provide monthly or yearly observations. These differences are
kept in mind during integration so that records with different meanings or time
levels are not incorrectly placed in the same relation.

The team also reviews entities that cannot be properly connected to the rest
of the model. Some collected datasets contain too many missing values, lack a
usable relationship key, or contain only a small amount of relevant
information. After discussion with the mentor, such entities are excluded
instead of creating relationships that are not supported by the source data.

The integration process reduces repeated entities and produces a smaller set
of combined relations. However, some of these relations still contain
repeated attributes, unsuitable keys, or different types of observations.
These problems are examined in the next section before the formal
normalization process begins.
\begin{figure}[H]
\centering

\begin{tikzpicture}[
    node distance=0.75cm and 1.2cm,
    source/.style={
        draw,
        rounded corners,
        align=center,
        minimum width=3.7cm,
        minimum height=0.8cm,
        font=\small
    },
    result/.style={
        draw=cublue,
        fill=cugrey,
        rounded corners,
        align=center,
        minimum width=5.2cm,
        minimum height=0.9cm,
        font=\small
    },
    arrow/.style={
        -{Stealth[length=2.5mm]},
        thick
    }
]

\node[source] (bmd) {Temperature Entity\\from BMD Data};
\node[source, right=of bmd] (brri) {Temperature Entity\\from BRRI Data};

\node[result, below=1.2cm of $(bmd)!0.5!(brri)$] (compare)
{Compare Meaning, Attributes, Keys,\\and Observation Level};

\node[result, below=of compare] (combined)
{Combined Temperature Record Relation};

\draw[arrow] (bmd) -- (compare);
\draw[arrow] (brri) -- (compare);
\draw[arrow] (compare) -- (combined);

\end{tikzpicture}

\caption{Example of entity comparison and integration}
\label{fig:entity-integration}
\end{figure}
Figure~\ref{fig:entity-integration} gives one example of the integration
process. Temperature entities identified from different sources are compared
before they are combined. The same process is followed for the other similar
environmental entities.

\subsection{Structural Problems and Model Decomposition}
\label{subsec:structural-decomposition}

After integrating the similar entities, the team reviews the structure of the
combined relations. At this stage, several problems remain. Some relations
contain different types of observations, some proposed primary keys contain
missing values, and some attributes do not depend on the same identifying
fields. These problems must be addressed before applying the normal-form
checks.

One major problem is found in the climate data. In the earlier model,
temperature, rainfall, humidity, wind, sunshine, radiation, and climatic-event
records are kept under a broad Climate Record entity. However, these
measurements do not have the same attributes or observation levels. For
example, some records are monthly, while sunshine and radiation records may
contain daily information. A single primary key is therefore not suitable for
all climate observations.

To solve this problem, the broad Climate Record entity is divided into
Temperature Record, Humidity Record, Rainfall Record, Wind Record, Sunshine
Record, Radiation Record, and Climatic Event Record. Each resulting relation
contains the attributes and key required for its own type of measurement.

The time information also requires revision. The earlier model uses one Time
entity for yearly, monthly, and daily records. This causes some records to
contain unnecessary or missing month and day values. The team therefore
creates separate Year Time, Month Time, and Day Time relations. Fiscal Year is
also kept separately because forest-area and industrial records use a start
year and an end year instead of one calendar date.

Similar changes are made to the industrial and establishment data. Industrial
Type is kept separately from Industry Usage, while Size is separated from Type
of Establishments. This prevents category names from being repeatedly stored
with observations from different fiscal years.

The main changes made during this stage are shown in
Table~\ref{tab:structural-revisions}. These changes prepare the relations for
the formal normalization process described in the next section.
\begin{figure}[H]
\centering

\begin{tikzpicture}[
    node distance=0.55cm,
    main/.style={
        draw=cublue,
        fill=cugrey,
        rounded corners=2pt,
        align=center,
        minimum width=5.2cm,
        minimum height=0.8cm,
        font=\small
    },
    process/.style={
        draw=cuteal,
        rounded corners=2pt,
        align=center,
        minimum width=5.2cm,
        minimum height=0.8cm,
        font=\small
    },
    relation/.style={
        draw,
        rounded corners=2pt,
        align=center,
        minimum width=3.4cm,
        minimum height=0.7cm,
        font=\small
    },
    arrow/.style={
        -{Stealth[length=2.4mm]},
        thick,
        draw=cuteal
    }
]

% Initial relation
\node[main] (climate) {Broad Climate Record};

% Decomposition rule
\node[process, below=0.6cm of climate] (separate)
{Separated by measurement type and time level};

% First row
\node[relation, below=0.85cm of separate, xshift=-3.8cm]
(temp) {Temperature Record};

\node[relation, right=0.35cm of temp]
(humidity) {Humidity Record};

\node[relation, right=0.35cm of humidity]
(rainfall) {Rainfall Record};

% Second row
\node[relation, below=0.35cm of temp]
(wind) {Wind Record};

\node[relation, right=0.35cm of wind]
(sunshine) {Sunshine Record};

\node[relation, right=0.35cm of sunshine]
(radiation) {Radiation Record};

% Third row
\node[relation, below=0.35cm of sunshine]
(event) {Climatic Event Record};

% Grouping box
\node[
    draw=cublue,
    dashed,
    rounded corners=3pt,
    fit=(temp)(humidity)(rainfall)(wind)(sunshine)(radiation)(event),
    inner sep=0.25cm
] (group) {};

% Only two arrows
\draw[arrow] (climate) -- (separate);
\draw[arrow] (separate) -- (group.north);

\end{tikzpicture}

\caption{Decomposition of the Climate Record entity}
\label{fig:climate-decomposition}
\end{figure}

Figure~\ref{fig:climate-decomposition} shows how the broad Climate Record
entity is divided according to measurement type. This allows each relation to
use a key and set of attributes that match its records.
Table~\ref{tab:structural-revisions} summarizes the problems found in the
integrated model and the changes made by the team.
\begin{table}[H]
\centering
\caption{Structural problems and revisions}
\label{tab:structural-revisions}

\small
\renewcommand{\arraystretch}{1.25}
\setlength{\tabcolsep}{6pt}

\begin{tabular}{@{}p{3.2cm}p{5.0cm}p{5.7cm}@{}}
\toprule

\textbf{Initial structure} &
\textbf{Problem identified} &
\textbf{Revision made} \\

\midrule

Broad Climate Record &
The measurements use different attributes, keys, and time levels. &
The entity is divided into Temperature, Humidity, Rainfall, Wind, Sunshine,
Radiation, and Climatic Event relations. \\

Unreliable candidate keys &
Some proposed key attributes contain missing values and cannot identify every
record. &
The affected keys are reviewed and replaced with suitable single or composite
keys. \\

Single Time entity &
Not every environmental record contains year, month, and day values. &
Year Time, Month Time, and Day Time are created as separate relations. \\

Calendar-year structure &
Forest and industrial records use a start year and an end year. &
A separate Fiscal Year relation is introduced for period-based records. \\

Combined industrial data &
Industry names are repeated with observations from different fiscal years. &
\texttt{Industry\_Type} is separated from \texttt{Industry\_Usage}. \\

Combined establishment data &
Size categories are repeated with period-based observations. &
Size is separated from Type of Establishments. \\

\bottomrule
\end{tabular}
\end{table}
After these revisions, the relations have clearer keys and contain records
with consistent meanings and observation levels. However, decomposition alone
does not prove that the relations are normalized. The relations are therefore
checked from 1NF to BCNF in the following sections.

\subsection{Theoretical Normalization Process}
\label{subsec:theoretical-normalization}

Normalization is used to organize data into relations with less repetition and
fewer data-dependency problems. The process examines how attributes depend on
the keys of a relation. When an unsuitable dependency is found, the relation
is divided into smaller relations without losing the connection between the
data \cite{codd1972further}.

The normal forms provide a sequence of checks. First Normal Form requires
atomic values and removes repeating groups. Second Normal Form checks whether
every non-key attribute depends on the complete primary key. Third Normal Form
removes dependencies between non-key attributes. Boyce--Codd Normal Form
applies a stricter condition by requiring every determinant to be a candidate
key \cite{kent1983five}.

The team applies these checks to the preliminary relations produced in the
previous sections. The relations are checked in order because a relation must
satisfy the earlier normal form before it can be examined for the next one.
Examples from the environmental datasets are used in the following sections
to show the changes made at each stage.
\begin{figure}[H]
\centering

\begin{tikzpicture}[
    node distance=0.5cm,
    stage/.style={
        draw,
        rounded corners=2pt,
        align=center,
        text width=9.5cm,
        minimum height=0.8cm,
        font=\small
    },
    first/.style={
        stage,
        draw=cublue,
        fill=cugrey
    },
    arrow/.style={
        -{Stealth[length=2.4mm]},
        thick,
        draw=cuteal
    }
]

\node[first] (initial)
{Preliminary relations identified from the source data};

\node[stage, below=of initial] (nf1)
{\textbf{First Normal Form (1NF)}\\
Remove repeating groups and keep one value in each field};

\node[stage, below=of nf1] (nf2)
{\textbf{Second Normal Form (2NF)}\\
Remove dependencies on only part of a composite key};

\node[stage, below=of nf2] (nf3)
{\textbf{Third Normal Form (3NF)}\\
Remove dependencies between non-key attributes};

\node[stage, below=of nf3] (bcnf)
{\textbf{Boyce--Codd Normal Form (BCNF)}\\
Check that every determinant is a candidate key};

\node[first, below=of bcnf] (final)
{Final normalized relations};

\draw[arrow] (initial) -- (nf1);
\draw[arrow] (nf1) -- (nf2);
\draw[arrow] (nf2) -- (nf3);
\draw[arrow] (nf3) -- (bcnf);
\draw[arrow] (bcnf) -- (final);

\end{tikzpicture}

\caption{Normalization process followed in the project}
\label{fig:normalization-process}
\end{figure}
Figure~\ref{fig:normalization-process} shows the order followed during
normalization. A relation moves to the next stage only after the requirements
of the previous normal form have been checked.
\subsection{First Normal Form}
\label{subsec:first-normal-form}

\subsubsection{Theoretical Requirement}

A relation is in First Normal Form when each field contains a single value and
there are no repeating groups \cite{kent1983five}. Each row must also represent
one identifiable record.

\subsubsection{Application to the Project}

Several source sheets contain daily or monthly values in separate columns. For
example, the sunshine dataset places the values for different days across the
same row. This format is suitable for presenting the data, but it is not
suitable for a relational table because the day columns form a repeating
group.

The team converts these columns into separate rows. The station, year, month,
and day are written explicitly for every observation. Table~\ref{tab:1nf-example}
shows a shortened example of this change.

\begin{table}[H]
\centering
\caption{Sunshine observations before and after First Normal Form}
\label{tab:1nf-example}
\small
\renewcommand{\arraystretch}{1.2}

\textbf{Before 1NF}\\[2mm]

\begin{tabular}{@{}lccccc@{}}
\toprule
\textbf{Station} &
\textbf{Year} &
\textbf{Month} &
\textbf{Day 1} &
\textbf{Day 2} &
\textbf{Day 3} \\
\midrule
Dhaka & 2019 & January & 6.1 & 5.9 & 6.0 \\
\bottomrule
\end{tabular}

\vspace{5mm}

\textbf{After 1NF}\\[2mm]

\begin{tabular}{@{}lccccl@{}}
\toprule
\textbf{Station} &
\textbf{Year} &
\textbf{Month} &
\textbf{Day} &
\textbf{Value} &
\textbf{Measure} \\
\midrule
Dhaka & 2019 & 1 & 1 & 6.1 & Sunshine Hours \\
Dhaka & 2019 & 1 & 2 & 5.9 & Sunshine Hours \\
Dhaka & 2019 & 1 & 3 & 6.0 & Sunshine Hours \\
\bottomrule
\end{tabular}
\end{table}

In the original structure, Day 1, Day 2, and Day 3 are repeated columns. After
1NF, the day number becomes an attribute, and every row represents one daily
observation. The same method is used for other source sheets containing
repeated day or month columns.

\subsection{Second Normal Form}
\label{subsec:second-normal-form}

\subsubsection{Theoretical Requirement}

A relation is in Second Normal Form when it is already in 1NF and every
non-key attribute depends on the complete primary key
\cite{kent1983five}. If an attribute depends on only part of a composite key,
a partial dependency exists.

\subsubsection{Application to the Project}

The 1NF Water Quality relation contains the station name, river name, water
category, year, parameter type, and measured value. One water-quality
observation is identified by the station, year, and parameter type.

The main dependency is:
\[
(\text{WQ Station Name},\text{Year},\text{Parameter Type})
\rightarrow \text{Value}
\]

However, the river name and water category depend only on the station:

\[
\text{WQ Station Name}
\rightarrow
\text{River Name},\text{Water Category}
\]

Therefore, the river information does not depend on the complete composite
key of the water-quality observation. The team separates the station-related
information from the measured values, as shown in
Table~\ref{tab:2nf-example}.

\begin{table}[H]
\centering
\caption{Water Quality relation before and after Second Normal Form}
\label{tab:2nf-example}
\small
\renewcommand{\arraystretch}{1.2}

\textbf{Before 2NF}\\[2mm]

\begin{tabular}{@{}p{2.8cm}p{2.5cm}p{1.7cm}cp{3.4cm}c@{}}
\toprule
\textbf{Station} &
\textbf{River} &
\textbf{Category} &
\textbf{Year} &
\textbf{Parameter} &
\textbf{Value} \\
\midrule
Mirpur Bridge &
Buriganga River &
Freshwater &
2018 &
Biochemical Oxygen Demand &
9.47 \\

Hajaribag &
Buriganga River &
Freshwater &
2018 &
Biochemical Oxygen Demand &
8.51 \\
\bottomrule
\end{tabular}

\vspace{5mm}

\textbf{After 2NF: River Station}\\[2mm]

\begin{tabular}{@{}p{4.2cm}p{3.8cm}p{2.8cm}@{}}
\toprule
\textbf{WQ Station Name} &
\textbf{River Name} &
\textbf{Water Category} \\
\midrule
Mirpur Bridge & Buriganga River & Freshwater \\
Hajaribag & Buriganga River & Freshwater \\
\bottomrule
\end{tabular}

\vspace{5mm}

\textbf{After 2NF: Water Quality}\\[2mm]

\begin{tabular}{@{}p{3.4cm}cp{4.2cm}c@{}}
\toprule
\textbf{WQ Station Name} &
\textbf{Year} &
\textbf{Parameter Type} &
\textbf{Value} \\
\midrule
Mirpur Bridge & 2018 & Biochemical Oxygen Demand & 9.47 \\
Hajaribag & 2018 & Biochemical Oxygen Demand & 8.51 \\
\bottomrule
\end{tabular}
\end{table}

After the decomposition, the Water Quality relation contains the measured
values, while River Station contains information that depends only on the
water-quality station name. This removes the partial dependency from the
measurement relation.

The source category remains in the 2NF and 3NF evidence, where it distinguishes
freshwater, lake, and marine material. It is not retained in the approved BCNF
relation because the final model represents the monitored water body through
\texttt{River\_Name}; the category adds no further identifying dependency.

\subsection{Third Normal Form}
\label{subsec:third-normal-form}

\subsubsection{Theoretical Requirement}

A relation is in Third Normal Form when it is in 2NF and a non-key attribute
does not determine another non-key attribute \cite{kent1983five}. Such a
dependency is known as a transitive dependency.

\subsubsection{Application to the Project}

A transitive dependency appears in the earlier groundwater-well structure. A
well identifies its district, while the district identifies its subdivision.
The subdivision therefore depends on the well indirectly through the district.

The dependencies can be written as:

\begin{gather*}
\text{Well ID}\rightarrow\text{District Name}\\
\text{District Name}\rightarrow\text{Subdivision}
\end{gather*}

Therefore:

\begin{gather*}
\text{Well ID}\rightarrow\text{District Name}\rightarrow\text{Subdivision}
\end{gather*}

To remove this dependency, the district and subdivision information is placed
in a separate relation. Table~\ref{tab:3nf-example} shows the structure before
and after that decomposition.

\begin{table}[htbp]
\centering
\caption{Groundwater-well structure before and after Third Normal Form}
\label{tab:3nf-example}
\small
\renewcommand{\arraystretch}{1.2}

\textbf{Before 3NF}\\[2mm]

\begin{tabular}{@{}lcllcc@{}}
\toprule
\textbf{Well ID} &
\textbf{Well No.} &
\textbf{District} &
\textbf{Subdivision} &
\textbf{Latitude} &
\textbf{Longitude} \\
\midrule
CM007 & 1 & Kumilla & Kumilla & 23.40 & 91.06 \\
CM002 & 2 & Kumilla & Kumilla & 23.58 & 91.15 \\
CM008 & 3 & Kumilla & Kumilla & 23.48 & 90.91 \\
\bottomrule
\end{tabular}

\vspace{5mm}

\textbf{After 3NF: Ground Water Well}\\[2mm]

\begin{tabular}{@{}lclcc@{}}
\toprule
\textbf{Well ID} &
\textbf{Well No.} &
\textbf{District} &
\textbf{Latitude} &
\textbf{Longitude} \\
\midrule
CM007 & 1 & Kumilla & 23.40 & 91.06 \\
CM002 & 2 & Kumilla & 23.58 & 91.15 \\
CM008 & 3 & Kumilla & 23.48 & 90.91 \\
\bottomrule
\end{tabular}

\vspace{5mm}

\textbf{After 3NF: District Subdivision}\\[2mm]

\begin{tabular}{@{}ll@{}}
\toprule
\textbf{District Name} &
\textbf{Subdivision} \\
\midrule
Kumilla & Kumilla \\
\bottomrule
\end{tabular}
\end{table}

After the decomposition, the well relation stores the district as a reference,
but the subdivision is stored only once for that district. The non-key
attribute Subdivision is therefore no longer repeated for every well.

The groundwater-well data are later excluded because the approved final ERD
has no groundwater entity. The records are nevertheless extracted and carried
through 3NF, so this example documents the transitive-dependency check performed
during schema refinement rather than a relation loaded into the final database.
\subsection{Boyce--Codd Normal Form}
\label{subsec:bcnf}

\subsubsection{Theoretical Requirement}

Boyce--Codd Normal Form applies a stricter rule than 3NF. A relation is in BCNF
when every determinant is a candidate key \cite{codd1972further}. A determinant
is an attribute, or a group of attributes, that determines another attribute.

\subsubsection{Application to the Project}

At this stage, the team reviews the determinants and candidate keys of the
retained relations. The key must match the level of the record. For example, a
temperature value is not identified only by station, year, and month because
the same station may report both minimum and maximum temperature in the same
month. Temperature Type must therefore be included in the candidate key.

Before the final key is selected, the following proposed determinant is
incomplete:
\[
(\text{Station Name},\text{Year},\text{Month})
\nrightarrow \text{Temperature}
\]
After including Temperature Type, the functional dependency becomes:
\[
(\text{Station Name},\text{Year},\text{Month},\text{Type})
\rightarrow \text{Temperature}
\]

\begin{table}[H]
\centering
\caption{Temperature Record candidate key established before the BCNF check}
\label{tab:bcnf-example}
\small
\renewcommand{\arraystretch}{1.2}

\textbf{Incomplete proposed key}\\[2mm]

\begin{tabular}{@{}lcccl@{}}
\toprule
\textbf{Station} &
\textbf{Year} &
\textbf{Month} &
\textbf{Temperature} &
\textbf{Meaning} \\
\midrule
Dinajpur & 2011 & 1 & 9.1  & Minimum \\
Dinajpur & 2011 & 1 & 20.8 & Maximum \\
\bottomrule
\end{tabular}

\vspace{2mm}

\begin{flushleft}
Proposed key:
\(\{\text{Station Name, Year, Month}\}\)
\end{flushleft}

\vspace{3mm}

\textbf{Corrected candidate key}\\[2mm]

\begin{tabular}{@{}lcccl@{}}
\toprule
\textbf{Station} &
\textbf{Year} &
\textbf{Month} &
\textbf{Type} &
\textbf{Temperature} \\
\midrule
Dinajpur & 2011 & 1 & Minimum & 9.1 \\
Dinajpur & 2011 & 1 & Maximum & 20.8 \\
\bottomrule
\end{tabular}

\vspace{2mm}

\begin{flushleft}
Final candidate key:
\(\{\text{Station Name, Year, Month, Type}\}\)
\end{flushleft}
\end{table}

Table~\ref{tab:bcnf-example} shows that station, year, and month cannot
determine one temperature value because minimum and maximum values can occur
for the same month. After Type is included, the complete determinant becomes
a candidate key. This corrects an incomplete proposed key; it is not presented
as a BCNF violation or decomposition. The BCNF test is applied only after the
candidate keys have been established.

The same check is carried out for the other final relations. Some examples are:
\begin{gather*}
(\text{Station Name},\text{Year},\text{Month})\rightarrow \text{Humidity}\\
(\text{Station Name},\text{Year},\text{Month})\rightarrow \text{Rainfall}\\
(\text{WQ Station Name},\text{Year},\text{Parameter Type})
\rightarrow \text{Value}\\
(\text{Industry Name},\text{Start Year},\text{End Year})
\rightarrow \text{Quantity},\text{Percentage}
\end{gather*}

In each case, the determinant is the candidate key of its relation. No further
decomposition is required at this stage because every non-trivial determinant
checked in the final relations is a candidate key. The final relations
therefore satisfy the BCNF requirement.

\subsection{Summary of Normalization}
\label{subsec:normalization-summary}

The normalization process begins by converting repeated spreadsheet columns
into atomic row-based records. Partial dependencies are then removed from
relations with composite keys, followed by the removal of transitive
dependencies. Finally, the determinants and candidate keys of the retained
relations are checked for BCNF. The resulting relations form the basis of the
final Entity--Relationship diagram presented in Section~\ref{sec:dm}.

Table~\ref{tab:normalization-summary} collects the worked examples of this
section into one view, so the structural change made at each stage can be read
without returning to the individual before-and-after tables.

\begin{table}[H]
\centering
\caption{Structural change introduced at each normalization stage}
\label{tab:normalization-summary}
\begin{tabular}{L{0.10\textwidth}L{0.25\textwidth}L{0.24\textwidth}L{0.26\textwidth}}
\toprule
Stage & Condition enforced & Project example & Structural outcome \\
\midrule
1NF & One value for each field and no repeating groups &
Sunshine values held in separate day columns &
Day columns become rows carrying station, year, month and day. \\
2NF & No non-key attribute depends on part of a composite key &
River name and water category depended only on the monitoring station &
Station description separates from the yearly measured values. \\
3NF & No non-key attribute depends on another non-key attribute &
Subdivision depended on district rather than on the well &
District detail is stored once and referenced by the well record. \\
BCNF & Every determinant is a candidate key &
Station, year and month did not identify one temperature value &
Temperature Type completes the candidate key before the BCNF test. \\
\bottomrule
\end{tabular}
\end{table}

\setlength{\abovedisplayskip}{11pt plus 3pt minus 6pt}
\setlength{\belowdisplayskip}{11pt plus 3pt minus 6pt}
\setlength{\abovedisplayshortskip}{0pt plus 3pt}
\setlength{\belowdisplayshortskip}{6.5pt plus 3.5pt minus 3pt}
\clearpage
